\section{Introduction}
\label{sec:intro}

Congruence closure is an important algorithm for satisfiability modulo theories
(SMT) solving~\cite{smt}, equality saturation~\cite{eqsat}, type
unification~\cite{unification}, and hardware verification~\cite{clausal-cc}.
Congruence closure determines whether a set of equalities and disequalities
between terms is satisfiable. For example, given the equalities $a = b$ and
$b = c$ and the disequality $f(a) \neq f(c)$, congruence closure will
conclude that this is unsatisfiable: $a = c$ follows by \emph{transitivity},
and then $f(a) = f(c)$ follows by \emph{congruence}, contradicting the
disequality. To decide this problem in practice, one must maintain
equalities using a union-find data structure and propagate congruences to
parent terms.


% Satisfiability modulo theories (SMT) solvers are used to decided the
% satisfiability of first-order logic formulas. State-of-the-art solvers such as
% Z3~\cite{z3} and cvc5~\cite{cvc5} can reason about a wide variety of
% ``theories'' such as uninterpreted functions, linear arithmetic, bit-vectors,
% and many others.

Congruence closure is a component of SAT solvers like Kissat~\cite{kissat};
equality saturation engines like egg~\cite{egg} and egglog~\cite{egglog}; and
SMT solvers like Z3~\cite{z3} and cvc5~\cite{cvc5}. These tools are used in a
variety of applications, including circuit equivalence
checking~\cite{clausal-cc}, software verification~\cite{dafny,verus}, hardware
design~\cite{hardware}, and mathematical theorem proving~\cite{lean-grind}.
Hence, their performance is critical, and much work has gone into improving
their efficiency. One promising avenue for improving the performance is
\emph{parallelization}: utilizing multiple processors to solve the problem
faster. 

To this end, we present two novel parallel congruence closure algorithms. Our
algorithms use a concurrent union-find data structure and a bulk-synchronous
parallel design to efficiently compute the congruence closure of a set of
equalities and disequalities.

% Sequential congruence closure algorithms are based on the seminal work by Nelson
% and Oppen~\cite{nelson-oppen-cc}, which maintains a worklist of pairs of terms
% that are known to be equal. The algorithm repeatedly processes pairs from the
% worklist, merging their equivalence classes in the union-find data structure and
% adding new pairs to the worklist until no more pairs can be added.

At a high level, our algorithms replace the classical sequential worklist of
Nelson and Oppen~\cite{nelson-oppen-cc} with a \emph{bulk-synchronous} closure
loop. Initial equalities are unioned into a lock-free concurrent
union-find~\cite{uf}. The algorithm then proceeds in rounds: each round
identifies a frontier of terms whose congruence class may have changed,
discovers all new congruences among them in parallel, and unions the resulting
classes. We implement each round using standard data-parallel primitives
provided by \textsc{ParlayLib}~\cite{parlaylib}. The algorithm terminates
when no round discovers a new merge, exactly mirroring the fixed point of
the sequential procedure.

The two algorithms differ in how they detect new congruences. The algorithm
\parentalgo maintains a parent list for each term, and when a term is involved
in a union, it adds all of its parents to a worklist. The second algorithm
\filteralgo marks terms as dirty when they are involved in a union and then
filters all terms to find those whose children are dirty.

 
In summary, this paper makes the following contributions:

\begin{itemize}
    \item We present two parallel congruence closure algorithms, \parentalgo
    and \filteralgo, both built on a bulk-synchronous parallel design that
    replaces the classical sequential worklist with rounds of data-parallel
    primitives over a lock-free concurrent union-find.
    \item We characterize when our algorithms achieve good speedup using the
    notions of congruence \emph{depth} and \emph{width}, and show empirically
    that width is the dominant factor.
    \item Despite congruence closure being P-complete (i.e. believed to not
    have a theoretically efficient parallel
    algorithm)~\cite{unification}, we demonstrate near-linear speedups up to
    32 cores over an efficient
    sequential baseline on random benchmarks, a synthetic benchmark suite
    designed to stress-test our algorithm, and a set of circuit equivalence
    benchmarks~\cite{clausal-cc}.
    \item We open source our \href{code and benchmarks}{https://github.com/amarshah10/ParallelEgraph}.
\end{itemize}
% We additionally discuss what would be required to integrate our algorithms
% into an existing SMT solver or equality saturation engine.

